% ============================================================
%  SCHOLA DOMESTICA — Magister Mathematica
%  Level 1 | Week 13 | Day 4
%  Topic: Oral Speed Drill #1 + Word Problems + Mixed Review
% ============================================================
\documentclass[12pt, letterpaper]{article}
\usepackage{sd_style}

\renewcommand{\lessonweek}{13}
\renewcommand{\lessonnum}{4}
\renewcommand{\lessontitle}{Oral Speed Drill \#1 \& Word Problems}

% IMAGE PROMPT (Beatrix Potter watercolour style):
% A small rabbit sits at a tiny wooden school desk under an oak tree, pencil in paw,
% ears alert, ready to begin. Autumn leaves drift past. White background, no text.

\begin{document}

\lessonheader{Week 13}{4}{Oral Speed Drill \#1 \& Word Problems}

{\small Name:\ansblanklg \hfill Date:\ansblank}
\goldrule

\begin{mdframed}[backgroundcolor=sd-gold!15,linecolor=sd-gold,linewidth=1.5pt,roundcorner=4pt,
  innerleftmargin=10pt,innerrightmargin=10pt,innertopmargin=8pt,innerbottommargin=8pt]
\textbf{\color{sd-gold}Oral Speed Drill \#1 — Teacher Script}\\[4pt]
\textit{Read each fact aloud at a steady pace. Child answers orally. Mark below.}\\
\textit{Goal: 20 of 25 correct. This is practice — celebrate effort!}\\[6pt]
\begin{multicols}{5}
\begin{enumerate}[leftmargin=*,label=\small\arabic*.,itemsep=1pt]
  \item $3-1$ \textit{[2]}
  \item $5-3$ \textit{[2]}
  \item $6-4$ \textit{[2]}
  \item $7-5$ \textit{[2]}
  \item $8-6$ \textit{[2]}
  \item $4-2$ \textit{[2]}
  \item $6-3$ \textit{[3]}
  \item $7-4$ \textit{[3]}
  \item $8-5$ \textit{[3]}
  \item $5-2$ \textit{[3]}
  \item $7-3$ \textit{[4]}
  \item $8-4$ \textit{[4]}
  \item $6-2$ \textit{[4]}
  \item $8-3$ \textit{[5]}
  \item $7-2$ \textit{[5]}
  \item $6-1$ \textit{[5]}
  \item $8-2$ \textit{[6]}
  \item $7-1$ \textit{[6]}
  \item $8-1$ \textit{[7]}
  \item $7-0$ \textit{[7]}
  \item $8-0$ \textit{[8]}
  \item $6-6$ \textit{[0]}
  \item $7-7$ \textit{[0]}
  \item $8-8$ \textit{[0]}
  \item $5-5$ \textit{[0]}
\end{enumerate}
\end{multicols}
\noindent Score:\underline{\hspace{1.5cm}}/25\hfill
\textit{Mastery: 20+. Below 15 — repeat drill tomorrow before advancing.}
\end{mdframed}

\bigskip
\noindent{\color{sd-blue}\large\textbf{Part I — Word Problems}} \hfill \textit{(Problems 1–6)}

\medskip\noindent\textit{Write a number sentence. Solve.}

\bigskip
\prob There are 8 apples in a basket. Tom eats 3. How many are left?\\[6pt]
\hspace{2em}Number sentence:\ansblanklg \quad Answer:\ansblank\ apples

\medskip
\prob Anna has 7 ribbons. She gives 4 to her sister. How many does Anna have now?\\[6pt]
\hspace{2em}Number sentence:\ansblanklg \quad Answer:\ansblank\ ribbons

\medskip
\prob Six birds sat on a fence. Five flew away. How many stayed?\\[6pt]
\hspace{2em}Number sentence:\ansblanklg \quad Answer:\ansblank\ bird(s)

\medskip
\prob Father caught 8 fish. He put 5 back. How many did he keep?\\[6pt]
\hspace{2em}Number sentence:\ansblanklg \quad Answer:\ansblank\ fish

\medskip
\prob A jar held 7 cookies. Some were eaten. Now there are 3 left. How many were eaten?\\[6pt]
\hspace{2em}Number sentence:\ansblanklg \quad Answer:\ansblank

\medskip
\prob Mother baked 8 rolls. Some were used at dinner. Now 2 remain. How many were used?\\[6pt]
\hspace{2em}Number sentence:\ansblanklg \quad Answer:\ansblank

\bigskip\bigskip
\noindent{\color{sd-blue}\large\textbf{Part II — Mixed Facts}} \hfill \textit{(Problems 7–16)}

\bigskip
\begin{multicols}{2}
\begin{enumerate}[leftmargin=*,label=\textbf{\arabic*.},itemsep=10pt,start=7]
  \item $4+4=\ansblank$
  \item $8-4=\ansblank$
  \item $3+5=\ansblank$
  \item $8-5=\ansblank$
  \item $2+6=\ansblank$
  \item $8-6=\ansblank$
  \item $1+7=\ansblank$
  \item $8-7=\ansblank$
  \item $5+3=\ansblank$
  \item $7-4=\ansblank$
\end{enumerate}
\end{multicols}

\bigskip
\begin{checkbox}
\textbf{\color{sd-green}Notice the Pattern!}\\[3pt]
Problems 7--8, 9--10, 11--12, 13--14 each use the \emph{same three numbers}.\\
$4+4=8 \;\longleftrightarrow\; 8-4=4$ \qquad $3+5=8 \;\longleftrightarrow\; 8-5=3$\\[3pt]
\textit{Knowing your addition facts makes subtraction fast!}
\end{checkbox}

\end{document}
